\clearpage\thispagestyle{empty}
\pdfbookmark[0]{Quellen zum Brief vom 3. August 1811}{lindenau1811aug-sources}
\begin{center}\textbf{Quellen zum Brief vom 3. August 1811}\end{center}
\noindent Diese KI-Übertragung ins heutige Deutsch folgt dem Brief von Gauß an Bernhard August von Lindenau in \href{https://gdz.sub.uni-goettingen.de/id/PPN236007467}{Gauß’ \textit{Werken}, Bd.~VI} (Göttingen, 1874), S.~326--327. Der Erstdruck erschien als Auszug in der \textit{Monatlichen Correspondenz}, Bd.~XXIV, August 1811, S.~\href{https://api.digitale-sammlungen.de/iiif/image/v2/bsb10538613_00190/full/1600,/0/default.jpg}{180}, \href{https://api.digitale-sammlungen.de/iiif/image/v2/bsb10538613_00191/full/1600,/0/default.jpg}{181} und \href{https://api.digitale-sammlungen.de/iiif/image/v2/bsb10538613_00192/full/1600,/0/default.jpg}{182}. Das digitalisierte Exemplar stammt aus der Bayerischen Staatsbibliothek.
\par\smallskip\noindent Die Zeitschrift bezeichnet Gauß als \textit{Ritter der westphälischen Krone}. In einer redaktionellen Fußnote auf S.~182 steht, dass die Ephemeride im nächsten Heft erscheinen werde. Die Werkausgabe lässt den Titel und die Fußnote weg und enthält die handschriftliche Bemerkung in eckigen Klammern, die im Erstdruck fehlt. KI-Übertragung und KI-Satz.
\clearpage
